\scanpage{12}
\handwritten{\mspage{3}}
Ceci posé, le cycle $Z$ (\unsure{ou|oa} plutôt, sa classe de cohomologie) définit un homomorphisme naturel
\[
Z:\quad R^1g_{\handwritten{*}}(\underline{Z}_{\unsure{\chi|\ell}}(1))\handwritten{\longrightarrow}R^{\unsure{2|a}n-1}f_{\handwritten{*}}(\underline{Z}_{\chi}(n))\ ,\tag*{(*)}
\]
d'où un homomorphisme
\[
Z:\quad H^1(S,R^1g_{\handwritten{*}}(\underline{Z}_{\chi}(1)))\handwritten{\longrightarrow}H^1(S,R^{2n-1}f_{\handwritten{*}}(\underline{Z}_{\chi}(n)))\ ,\tag*{(**)}
\]
et la compatibilité générale énoncée plus haut nous dit que $u(x)=u(c\chi(z))$ n'est autre que l'image de $u(x')$ du premier membre de \deleted{\unsure{(*)x|(*)}} (**). Nous \unsure{aurons|avrons} besoin simplement pour la suite de savoir que \uline{si l'homo\deleted{\unsure{morphismo|morphisme}}morphisme\deleted{xx}} (*) \uline{est nul} (\uline{ou ce qui revient au même, si l'homomorphisme correspondant sur les fibres géométriques en} $\bar s$ \deleted{est}
\[
H^1(C_{\bar s},\underline{Z}_{\chi}(1))\handwritten{\longrightarrow}H^{2n-1}(Y_{\bar s},\underline{Z}_{\chi}(n))\tag*{(* bis)}
\]
\uline{est nulle}), \uline{alors il en est de même de} $u(x)\handwritten{\in}H^1(S,R^{2n-1}f_{\handwritten{*}}(\underline{Z}_{\chi}(n)))$.

2.\quad Soit
\[
i:Y\handwritten{\longrightarrow}X
\]
un morphisme de schémas pro\deleted{jectif}\handwritten{pre}s et lisses sur un corps algébriquement clos $k$, $n$ un entier. On \deleted{suppose} considère l'homomorphisme
\[
i^{\handwritten{*}}:\quad H^{2n-1}(X,\underline{\deleted{Z}\handwritten{Q}}_{\chi}(n))\handwritten{\longrightarrow}H^{2n-1}(Y,\underline{\deleted{Z}\handwritten{Q}}_{\chi}(n))\ ,\tag*{\handwritten{(***)}}
\]
et on suppose:

a)\quad Il existe un homomorphisme en sens inverse $\unsure{v*|v^*|vk}$, dont la restriction à l'image de\deleted{jxx} $i^{\handwritten{*}}$ soit un inverse à droite de $i^{\handwritten{*}}$, et qui soit ``algébrique'' i.e. donné par un cycle algébrique sur $Y\mathbin{\mathrm{x}}X$. (NB Cette hypothèse est un cas particulier des ``conjectures standart'', donc est sans doute touhours vérifiée; elle est vérifiée en tous cas si $i$ est l'i\deleted{\unsure{x}}nclusion d'une section hyperplane dans $X$, lorsque l'opération $\handwritten{\Lambda}$ sur \deleted{X est} la cohomologie de $X$ est algébrique, \deleted{ce} \inserted{propriété} qu'on sait vérifier dans de très nombreux cas, par exemple lorsque $X$ se remonte en car. zéro avec sa polarisation.)
% archive sheet number: 11
